% 02-project-planning. The repository's own three-step vocabulary — research,
% then plan, then implement (merra_modis_comparison/CLAUDE.md, "Workflow") —
% drawn as a chain, with a review gate between phases and the artifact each
% phase left behind counted under it.
%
% The counts are real and were measured, not estimated:
%   research/  3 notes   FSCA_PRODUCT_NOTES, SNOTEL_PRODUCT_NOTES, COARSE_DEM_NOTES
%   plan/      2 specs   SNOWPACK_REANALYSIS_PLAN (210 ln), SEASON_SHAPE_PLAN (188 ln)
%              (the two archived FSCA plans beside them describe the superseded
%              experiment and are deliberately not counted)
%   src/      26 modules (excluding __init__.py)
%   tests/   329 tests across 14 files
% Re-measure before the talk if the branch has moved.
%
% The two boxes are the point of the slide. A plan that only fixed WHAT TO
% COMPUTE would be unremarkable; this one also fixed WHAT MAY BE CLAIMED from
% the result, and that half is what kept a plausible wrong headline off the
% figures. Both lists are compressed from CLAUDE.md's own sections ("Scientific
% invariants" and "Claims discipline").
\begin{frame}[t]
\frametitle{Project Planning}
\vspace{0.30em}
% The chain. Phases in the accent, gates in NASA red — the review between
% phases is the only thing on this slide the palette's accent colour is spent
% on, because it is the mechanism the slide is about.
\begin{center}
\begin{tikzpicture}[
  ph/.style={draw=accent, very thick, rounded corners=4pt, fill=accentFill,
             minimum width=2.45cm, minimum height=0.84cm,
             font=\small, text=accentDeep, align=center},
  art/.style={font=\scriptsize, text=mutedInk, align=center, anchor=north},
  gate/.style={draw=accentSoft, very thick, fill=white, circle, inner sep=0.5pt,
               minimum size=0.46cm, font=\scriptsize, text=accentSoft},
  fl/.style={draw=accent, very thick}]

  \node[ph] (p1) at (0.00,0) {Research};
  \node[ph] (p2) at (3.85,0) {Plan};
  \node[ph] (p3) at (7.70,0) {Implement};
  \node[ph] (p4) at (11.55,0) {Verify};

  \foreach \a/\b/\x in {p1/p2/1.925, p2/p3/5.775, p3/p4/9.625} {
    \draw[fl] (\a) -- (\b);
    \node[gate] at (\x,0) {$\checkmark$};
  }

  % What each phase left on disk. Counted, not characterised.
  \node[art] at (0.00,-0.56)  {\codepath{research/}\\3 product notes};
  \node[art] at (3.85,-0.56)  {\codepath{plan/}\\2 specifications};
  \node[art] at (7.70,-0.56)  {\codepath{src/}\\26 modules};
  \node[art] at (11.55,-0.56) {\codepath{tests/}\\329 tests};
\end{tikzpicture}
\end{center}
\vspace{0.35em}
% TITLES ARE A MATCHED PAIR — "Fixed Before Any Code" / "Fixed Before Any
% Number" — and each stands on its own. Two concrete nouns, so they are
% parallel, and the pair carries the argument the slide exists for: the timing.
% Before a single number was quoted, the plan had settled which numbers the
% data would be able to support. That is literally true, not a flourish — the
% right-hand set was established by measurement over WY2000-WY2026 before the
% full record was built (plan/SNOWPACK_REANALYSIS_PLAN.md, section 6.1).
%
% TWO EARLIER VERSIONS FAILED. Do not go back to either.
%
% "The Plan Fixes" / "And What May Be Claimed" was written to be read across
% the gap as one sentence. A box title is a label, not a clause fragment, and
% nobody assembles a sentence across two separately framed boxes — least of all
% two in different colours, which says outright that they are different things.
% The join did not even parse: the left box takes noun phrases and completes as
% a sentence; the right-hand bullets were full clauses.
%
% "Fixed Before Any Claim" with imperative bullets then failed on VOCABULARY,
% which is the harder fault because it builds clean and reads fluently. Three
% terms said nothing to the room:
%
%   "every number names its product"  — devoid of meaning. "Product" is archive
%     vocabulary for a dataset, and the sentence never says WHY. The real point
%     is that the two models disagree about amount by a factor of three or
%     more, so any snow amount is a fact about THAT MODEL and not about
%     Colorado. The bullet now says which model, in those words.
%   "load-bearing"  — an engineering metaphor. The real point is that the last
%     months of WY2026 come from ERA5T, the preliminary stream, which can still
%     be revised. The bullet now says preliminary.
%   "claim"  — reads as jargon in a title. The box is about what the result may
%     say, and "number" is the concrete thing every bullet is actually about.
%
% Bullets are noun phrases, matching the left box, so both lists genuinely
% complete their own title. Each carries its REASON where it has room —
% "not MERRA-2 in April, it has melted out" is a rule and its justification in
% seven words, which the bare rule was not. Sources, in order:
% CLAUDE.md "Claims discipline" bullets 5, 5, 2, 3 and 6.
%
% Every term here should survive being said out loud to someone who does not
% work on this. That is the test the previous two versions failed.
\begin{columns}[T,onlytextwidth]
\begin{column}{0.395\textwidth}
\begin{block}{Fixed Before Any Code}
\small
\begin{itemize}\itemsep0.12em
  \item the question
  \item domain and period
  \item variables, units, conversions
  \item the weighting rule
  \item what is out of scope
\end{itemize}
\end{block}
\end{column}
\begin{column}{0.575\textwidth}
% Alerted, so the pair reads as constraint against specification rather than as
% two lists of the same kind. Under anthropic_orange that is the dark brown —
% the grave tier, not a louder one; see themes/anthropic_orange.tex.
\begin{alertblock}{Fixed Before Any Number}
\small
\begin{itemize}\itemsep0.12em
  \item which years were extreme, not how much snow
  \item which model each number came from
  \item daily data for peak snowpack, never weekly
  \item not MERRA-2 in April, it has melted out
  \item where WY2026 rests on preliminary ERA5T
\end{itemize}
\end{alertblock}
\end{column}
\end{columns}
\end{frame}
